% !TEX TS-program = pdflatexmk
\documentclass[12pt]{amsart}
\usepackage{multicol}
\usepackage{float}

%\usepackage[parfill]{parskip}    % Activate to begin paragraphs with an empty line rather than an indent

\usepackage[margin=1in]{geometry}


\usepackage{amsmath,amssymb,amsthm,latexsym,graphicx}
\usepackage[normalem]{ulem}
\usepackage{setspace} %used for doublespacing, etc.
\usepackage{hyperref}
\usepackage{cancel}
\usepackage[dvipsnames,usenames]{color}
\usepackage[all]{xy}
\usepackage{fancyhdr}
\pagestyle{fancy}
	\renewcommand{\headrulewidth}{0.5pt} % and the line
	\headsep=1cm
	
\DeclareGraphicsRule{.tif}{png}{.png}{`convert #1 `dirname #1`/`basename #1 .tif`.png}

%Some useful environments.
\newtheorem{theorem}{Theorem}
\newtheorem{corollary}[theorem]{Corollary}
\newtheorem{conjecture}[theorem]{Conjecture}
\newtheorem{lemma}[theorem]{Lemma}
\newtheorem{proposition}[theorem]{Proposition}
\newtheorem{definition}[theorem]{Definition}
\newtheorem{example}[theorem]{Example}
\newtheorem{axiom}{Axiom}
\theoremstyle{remark}
\newtheorem{remark}{Remark}
\newtheorem*{exercise}{Exercise}%[section]

%Some useful shortcuts for our favorite fields
\def\RR{\ensuremath{\mathbb R}} %Note, you can use these WITHOUT entering math mode
\def\NN{\ensuremath{\mathbb N}}
\def\ZZ{\ensuremath{\mathbb Z}}
\def\QQ{{\ensuremath\mathbb Q}}
\def\CC{\ensuremath{\mathbb C}}
\def\EE{{\ensuremath\mathbb E}}

%Some useful shortcuts for formatting lists
\newcommand{\bc}{\begin{center}}
\newcommand{\ec}{\end{center}}
\newcommand{\be}{\begin{enumerate}}
\newcommand{\ee}{\end{enumerate}}
\newcommand{\bi}{\begin{itemize}}
\newcommand{\ei}{\end{itemize}}

%Some useful shortcuts for formatting mathematical symbols
\newcommand{\ol}[1]{\overline{#1}}
\newcommand{\oimp}[1]{\overset{#1}{\iff}} %labeled iff symbol
\newcommand{\bv}[1]{\ensuremath{ \vec{\mathbf{#1}}} } %makes a vector.
\newcommand{\mc}[1]{\ensuremath{\mathcal{#1}}} %put something in caligraphic font
\newcommand{\mpg}[1]{\marginpar{ #1}} %to put comments in margins
\newcommand{\bsl}[1]{\texttt{\symbol{92}{\em #1}}} %for backslashes.
\newcommand{\normale}{\trianglelefteq}
\newcommand{\normal}{\triangleleft}

%Code for formatting the proofs a little nicer for submitted homework
\makeatletter
\renewenvironment{proof}[1][\proofname]{\par\doublespacing
  \pushQED{\qed}%
  \normalfont \topsep6\p@\@plus6\p@\relax
  \list{}{%
    \settowidth{\leftmargin}{\itshape\proofname:\hskip\labelsep}%
    \setlength{\labelwidth}{0pt}%
    \setlength{\itemindent}{-\leftmargin}%
  }%
  \item[\hskip\labelsep\itshape#1\@addpunct{:}]\ignorespaces
}{%
  \popQED\endlist\@endpefalse
  \singlespacing
}
\makeatother


%Commenting tools. You can ignore this stuff unless we're exchanging materials where I'm commenting in detail on how you are writing/using LaTeX.
\usepackage{soul}
\definecolor{highlight}{rgb}{1,0.6,0.6}
\sethlcolor{highlight}
\newcommand{\hlm}[1]{\colorbox{highlight}{$\displaystyle #1$}}
\newtheoremstyle{mycomment}{\topsep}{-0in}{\small \itshape \sffamily}{}{\small \itshape\sffamily}{:}{.5em}{}
\theoremstyle{mycomment}
\newtheorem*{acomment}{\color{BrickRed}{Comment}}
\newcommand{\com}[1]{{\color{OliveGreen}\begin{acomment}{#1} %#2 \color{black} 
\end{acomment}\noindent}}
%\newcommand{\com}[1]{{\color{BrickRed}{\\ Comment:}\color{OliveGreen}{#1} \\}}
\newcommand{\red}[1]{{\color{BrickRed} #1}}
\newcommand{\blue}[1]{{\color{MidnightBlue}#1}}
\newcommand{\green}[1]{{\color{OliveGreen}#1}}
\newcommand{\mwrong}[2]{\red{\cancel{#1}}\green{#2}}
\newcommand{\wrong}[2]{\red{\sout{#1}}\green{#2}}
\definecolor{OliveGreen}{rgb}{.3,.5,.2}
\definecolor{MidnightBlue}{rgb}{.3,.4,.6}
\newcommand{\pts}[1]{\hfill\blue{{#1}/5}}

\chead{MATH 631F}
\pagestyle{fancy}
%Modify these items:
\rhead{\emph{Algebra Group}}
\lhead{\emph{HW \#8 --- \today}}

\DeclareMathOperator{\lcm}{lcm}
\begin{document}

\thispagestyle{fancy}
\author{Coordinator: Glen Woodworth}


%\usepackage{float} <- Throw this in the preamble



\begin{exercise}[\textbf{5.4.17} --- Glen Woodworth] If $K$ is normal subgroup of $G$ and $K$ is cyclic, prove that $G'\le C_G(K)$. [Recall that the automorphism group of a cyclic group is abelian.]
  \begin{proof} Let $K\trianglelefteq G$. Recall that the automorphism group of a cyclic group is abelian. Since $K$ is cyclic, Aut$(K)$ is abelian. Consider the map
  \begin{align*}
  \Phi:G&\rightarrow \text{Aut }(K)\\
  g&\mapsto\phi_g:k\mapsto gkg^{-1}.\end{align*}
  Since $K\trianglelefteq G$, $gkg^{-1}\in K$ for all $k\in K$ and for all $g\in G$. Thus, $\Phi$ is well-defined. We now show that $\Phi$ is a homomorphism. Let $g_1,g_2\in G$ and $k\in K$. Observe,
  \begin{align*}
  \Phi(g_1g_2)(k)&=\phi_{g_1g_2}(k)=g_1g_2k(g_1g_2)^{-1}=g_1g_2kg_2^{-1}g_1^{-1}=\phi_{g_1}(g_2kg_2^{-1})\\
  &=\phi_{g_1}(\phi_{g_2}(k))=(\phi_{g_1}\circ\phi_{g_2})(k)=(\Phi(g_1)\circ\Phi(g_2))(k).
  \end{align*}
Thus, $\Phi$ is a homomorphism. We now show that $\ker\Phi= C_G(K)$. Let $x\in C_G(K)$. Then \[\Phi(x)(k)=xkx^{-1}=k.\]
Thus, $\Phi(x)=Id$. Suppose $y\in\ker\Phi$. We have \[k=\Phi(y)(k)=yky^{-1}.\] Therefore, $y\in C_G(K)$. Hence, $\ker\Phi=C_G(K)$. Proposition 7 from Section 5.4 states that if $\phi:G\rightarrow A$ is any homomorphism where $G$ is a group and $A$ is an abelian group, then $G'\le\ker\phi$. Since Aut$(K)$ is abelian, and $\Phi$ is a homomorphism with $\ker\Phi=C_G(K)$, we have $G'\le\ker\Phi=C_G(K)$.
  \end{proof}
\end{exercise}



\begin{exercise}[\textbf{5.5.8} --- Glen Woodworth] Construct a non-abelian group of order $75$. Classify all groups of order 75 (there are three of them). [Use Exercise 6 to show that the non-abelian group is unique.] (The classification of groups of order $pq^2$, where $p$ and $q$ are primes with $p<q$ and $p$ not dividing $q-1$, is quite similar.)
  \begin{proof} Consider the homomorphism \begin{align*}\Phi:Z_{25}&\rightarrow\text{Aut}(Z_{3})\\
\overline{0}&\mapsto\text{Id}:\overline{1}\mapsto\overline{1}\\
\overline{1}&\mapsto \phi:\overline{1}\mapsto \overline{2}.
  \end{align*} 
  We show that $\Phi$ is well-defined. Observe,\[\phi(\overline{1}+\overline{2})=\phi(\overline{0})=\overline{0}=\overline{2}+\overline{1}=\phi(\overline{1})+\phi(\overline{2})\]
Thus, $\phi$ is a homomorphism. Since $\phi$ is obviously a bijection, $\phi\in$Aut$(Z_3)$. Thus, $\Phi$ is well-defined. We now show $\Phi$ is a homomorphism. First, observe that since $\Phi(\overline{1})(\overline{1})=\phi(\overline{1})=\overline{2}$ and since there are only two options in the image, if $x\in Z_{25}$ is odd, then $\Phi(x)(\overline{1})=\phi(\overline{1})=\overline{2}$. Similarly, if $y\in Z_{25}$ is even, then $\Phi(y)(\overline{1})=$Id$(\overline{1})=\overline{1}$. Now, suppose $g_1,g_2\in Z_{25}$. We have three cases to consider. For the first case, suppose $g_1$ and $g_2$ are both even. Then $g_1+g_2$ is even. Observe,\[
\Phi(g_1+g_2)(\overline{1})=\text{Id}(\overline{1})=\text{Id}(\text{Id}(\overline{1}))=\Phi(g_1)(\Phi(g_2)(\overline{1}))=(\Phi(g_1)\circ\Phi(g_2))(\overline{1}).\]
For the second case, suppose $g_1$ and $g_2$ are both odd. Then $g_1+g_2$ is even. We have
\[\Phi(g_1+g_2)(\overline{1})=\text{Id}(\overline{1})=\overline{1}=\phi(\overline{2})=\phi(\phi(\overline{1}))=\Phi(g_1)(\Phi(g_2)(\overline{1}))=(\Phi(g_1)\circ\Phi(g_2))(\overline{1}).\]
For the last case, suppose $g_1$ is even and $g_2$ is odd. Then $g_1+g_2$ is odd. Observe,
\[\Phi(g_1+g_2)(\overline{1})=\phi(\overline{1})=\text{Id}(\phi(\overline{1}))=\Phi(g_1)(\Phi(g_2)(\overline{1}))=(\Phi(g_1)\circ\Phi(g_2))(\overline{1}).\]
Thus, $\Phi$ is a well-defined homomorphism. Hence, by the definition of semidirect product and Theorem 10, $G=Z_3\rtimes_\Phi Z_{25}$ is a group of order $75$. Let $\cdot$ denote the left action of $Z_{25}$ on $Z_3$ determined by $\Phi.$ To see that $G$ is nonabelian, consider $(\overline{1},\overline{3}),(\overline{2},\overline{4})\in G$. Observe,
\begin{align*}
(\overline{1},\overline{3})(\overline{2},\overline{4})(\overline{1},\overline{3})^{-1}&=(\overline{1}+\overline{3}\cdot\overline{2},\overline{3}+\overline{4})(\overline{22}\cdot\overline{2},\overline{22})=(\overline{1}+\Phi(\overline{3})(\overline{2}),\overline{7})(\Phi(\overline{22})(\overline{2}),\overline{22})\\
&=(\overline{1}+\phi(\overline{2}),\overline{7})(\text{Id}(\overline{2}),\overline{22})=(\overline{1}+\overline{1},\overline{7})(\overline{2},\overline{22})=(\overline{2},\overline{7})(\overline{2},\overline{22})\\
&=(\overline{2}+\overline{7}\cdot\overline{2},\overline{7}+\overline{22})=(\overline{2}+\Phi(\overline{7})(\overline{2}),\overline{4})=(\overline{2}+\phi(\overline{2}),\overline{4})=(\overline{2}+\overline{1},\overline{4})\\
&=(\overline{0},\overline{4}).
\end{align*}
Since $(\overline{2},\overline{4})\ne(\overline{0},\overline{4})$, $G$ is nonabelian. We are given that there are only three groups of order 75. The other two groups of order 75 are $Z_{75}$ and $Z_{15}\times Z_{5}$, which are both abelian. Note that since $(5,15)\ne 1$, $Z_{75}$ is not isomorphic to $Z_{15}\times Z_{5}$. Exercise 6 of section 5.5 states that if there exists another homomorphism $\Phi_1$ from $Z_{25}$ to Aut$(Z_3)$ such that $\Phi(Z_{25})$ and $\Phi_1(Z_{25})$ are conjugate subgroups of Aut$(Z_3)$, then $Z_3\rtimes_{\Phi}Z_{25}\cong Z_3\rtimes_{\Phi_1}Z_{25}.$ If we suppose that $\Phi_1$ is non-trivial, then $\phi\in\Phi_1(Z_{25})$ because $\phi$ is the only non-identity element in Aut$(Z_3)$. Since Id$\in\Phi_1(Z_{25})$ by the definition of a homomorphism, $\Phi_1(Z_{25})=\Phi(Z_{25})$. By Exercise 6, $Z_3\rtimes_{\Phi}Z_{25}\cong Z_3\rtimes_{\Phi_1}Z_{25}.$
  \end{proof}
 
\end{exercise}

\begin{exercise}[\textbf{5.1.5} --- Glen Woodworth] Exhibit a nonnormal subgroup of $Q_8\times Z_4$ (note that every subgroup of each factor is normal).
  \begin{proof} Consider the subgroup $\langle (-k,0)\rangle=\{(-k,0),(1,0)\}$ in $Q_8\times Z_4$. Observe, \begin{align*}
  (j,0)(-k,0)(-j,0)&=(j(-k),0)(-j,0)=(-jk,0)(-j,0)=(i,0)(-j,0)\\
  &=(i(-j),0)=(-ij,0)=(k,0)\notin\langle (-k,0)\rangle.\end{align*}
 Hence, $\langle (-k,0)\rangle$ is not a normal subgroup of $Q_8\times Z_4$.
  \end{proof}
\end{exercise}




  \end{document}